% chapter3_system_study.tex
% Automated Diagram Digitization -- System Study Chapter

\chapter{System Study}
\label{ch:systemstudy}

The system study phase involves understanding existing diagram 
digitization challenges and analyzing requirements for the Automated 
Diagram Digitization system. This chapter covers problem definition, 
feasibility analysis, and requirements gathering. A comprehensive 
investigation was conducted through literature review, analysis of 
existing diagramming tools, and identification of user needs to 
establish a solid foundation for system design and implementation.

%---------------------------------------------------------------
\section{Existing System}
\label{sec:existing-system}
%---------------------------------------------------------------

The current landscape of diagram digitization tools presents several 
significant challenges that limit their effectiveness in real-world 
usage. This section examines the key problems encountered by users 
today and the disadvantages that motivate the development of a more 
intelligent approach.

% C3: Merged Sections 3.1.1 and 3.1.2 into one unified subsection
\subsection{Challenges and Disadvantages of Existing Systems}
\label{sec:existing-challenges}

The existing process of digitizing flowcharts and diagrams faces several critical challenges that impact user productivity and diagram quality. These difficulties arise from the fundamental limitations of current tools, which were designed primarily for manual creation rather than automated recognition. One of the primary obstacles is that manual redrawing is almost always required, with users forced to recreate diagrams using tools such as Microsoft Visio or Lucidchart—a time-consuming process that can take up to 30 minutes for moderately complex flowcharts and introduces transcription errors. This inefficiency is compounded by a total lack of automated shape recognition, meaning current systems cannot automatically identify shapes, text, or connections from images of hand-drawn diagrams, requiring each element to be identified and entered individually.

Furthermore, standard OCR tools often fail to preserve text formatting and line breaks, resulting in unreadable merged text that requires extensive manual correction. A significant connection detection gap also exists, where existing systems fail to recognize directional relationships, leaving researchers and developers to manually draw arrows and define flow logic. Most existing tools are further limited by restricted input methods, typically supporting only manual drag-and-drop rather than image uploads or text-to-flowchart generation. Beyond these issues, the absence of intelligent editing features—such as auto-layout and undo/redo capabilities—and restricted export options (often limited to raster images) further hinder the usability of digitized outputs across professional workflows. Finally, a lack of theme or accessibility support reduces comfort during extended usage, collectively making the current landscape of diagram digitization highly suboptimal.

%---------------------------------------------------------------
\section{Proposed System}
\label{sec:proposed-system}
%---------------------------------------------------------------

The Automated Diagram Digitization system is designed to address all 
of the shortcomings identified in the existing system analysis. This 
section provides a high-level overview of the proposed solution and 
outlines the specific advantages it offers over traditional 
approaches.

\subsection{System Overview}
\label{sec:system-overview}

The Automated Diagram Digitization system is a web-based platform that combines AI-powered recognition with an interactive diagram editor to provide a complete workflow from printed images to editable digital diagrams. Central to the system is its AI-powered recognition module, which leverages Google’s Gemini AI with custom prompts to analyze flowchart images, identify shapes such as processes and decisions, and extract embedded text with connection paths. To accommodate diverse user workflows, the platform supports multiple input methods, including image uploads, natural language descriptions for algorithm-to-flowchart conversion, and direct JSON import.

The interactive editor module serves as the core workspace, providing professional-grade features such as drag-and-drop from a shape palette, snapshot-based undo/redo, node resizing with text adaptation, and automatic layout alignment using the Dagre engine. Users can further refine their diagrams through extensive text customization, including bold, italic, and font size adjustments. Once finalized, diagrams can be exported in multiple formats, such as PNG, SVG, and JSON, ensuring both accessibility and portability. The system also prioritizes user experience with theme support for dark and light modes and ensures workflow continuity through persistent storage in a PostgreSQL database with UUID-based identification.

\subsection{Advantages of Proposed System}
\label{sec:advantages}

The following advantages illustrate how the proposed system directly 
addresses the limitations identified in the existing system analysis, 
delivering a more automated, flexible, and feature-rich experience 
for end users.
The proposed system offers several key benefits that enhance both 
usability and efficiency. Automatic shape identification significantly 
reduces manual redrawing, thereby saving considerable time and effort 
per diagram. Text extraction with proper line break preservation 
maintains readability and retains the original formatting of the 
source diagram. Automatic connection detection ensures that edges are 
created with correct directional arrows, preserving the logical flow 
of the diagram.

Furthermore, the availability of multiple input methods provides 
flexibility to accommodate diverse user workflows and use cases. 
Comprehensive editing features, comparable to those of professional 
diagramming tools, enable effective post-recognition refinement. 
Multiple export options support a wide range of applications, including 
documentation, data portability, and system integration.

In addition, the inclusion of undo and redo functionality with snapshot 
history prevents accidental data loss and allows users to experiment 
safely during editing. Auto-layout functionality contributes to a 
professional appearance by ensuring proper alignment of diagram elements 
after recognition or manual adjustments. Theme switching between dark 
and light modes enhances user comfort by reducing eye strain during 
extended usage. Finally, persistent storage capabilities allow users to 
maintain workflow continuity across multiple sessions.

%---------------------------------------------------------------
\section{Feasibility Study}
\label{sec:feasibility}
%---------------------------------------------------------------

A thorough feasibility study was conducted to evaluate whether the 
proposed system can be successfully developed and deployed within the 
available resources. The study examines technical, operational, and 
economic dimensions to confirm that the project is both achievable 
and sustainable.

\subsection{Technical Feasibility}
\label{sec:tech-feasibility}

The project utilizes proven, mature technologies that have demonstrated reliability across numerous deployments. Each component was selected based on its stability, community support, and suitability for the specific requirements of the platform. The backend is powered by FastAPI, which provides high-performance asynchronous API endpoints and automatic documentation, facilitating efficient development and testing. AI integration is achieved through the Google Gemini API, offering state-of-the-art multimodal capabilities for image analysis and text generation. On the frontend, React 18 with Vite ensures optimal performance and efficient rendering, while React Flow provides specialized diagramming capabilities. Data persistence is managed by PostgreSQL with JSONB support, ensuring ACID compliance and a flexible schema for complex flowchart data. Finally, the system’s modular architecture supports cloud deployment with horizontal scaling capabilities as user numbers grow.

Each technology in the stack has extensive community support and a 
demonstrated track record in application development, collectively 
reducing implementation risk.

\subsection{Operational Feasibility}
\label{sec:op-feasibility}

Beyond technical considerations, the system is designed to be practically usable by the intended audience with minimal training or specialized knowledge. The platform features a user-friendly interface with a tabbed layout that clearly presents options for image upload, text prompts, and JSON import. Interaction within the editor is intuitive, incorporating drag-and-drop from a shape palette, familiar keyboard shortcuts like Ctrl+Z and Ctrl+Y, and double-click edge labeling that follows common software conventions. To ensure clarity during usage, the system provides clear visual feedback on upload progress, processing status, and errors. Furthermore, the system is architected for minimal maintenance through automatic error recovery and fallback mechanisms. Responsive design principles ensure that the interface remains usable across a wide range of devices, including desktops, tablets, and mobile phones.

\subsection{Economic Feasibility}
\label{sec:economic-feasibility}

The economic feasibility of the proposed system was assessed by examining development, deployment, and maintenance costs, confirming that the project can be delivered within a reasonable budget while providing measurable productivity gains. By utilizing open-source technologies such as FastAPI, React, and PostgreSQL, the initial development investment is minimized. The Google Gemini API further enhances cost-effectiveness by offering competitive pricing and free-tier options suitable for the development phase. The automation of diagram digitization directly translates to reduced labor costs by eliminating hours of manual redrawing effort. Additionally, the system’s scalable cloud architecture allows for a pay-as-you-grow approach, avoiding large upfront infrastructure costs. Ultimately, educational institutions and businesses can achieve a positive return on investment through improved efficiency and significantly reduced manual digitization effort.

%---------------------------------------------------------------
\section{Requirements Analysis}
\label{sec:requirements}
%---------------------------------------------------------------

Requirements analysis is a critical phase that bridges the gap 
between problem identification and system design. It involves 
identifying, documenting, and organizing the specific capabilities 
and constraints that the system must fulfill. The requirements for 
the Automated Diagram Digitization system were gathered through 
analysis of existing tools, evaluation of user needs, and review of 
the challenges identified earlier in this chapter.

\subsection{Functional Requirements}
\label{sec:functional-req}

Based on user needs and system objectives, the following functional 
requirements have been identified. These requirements define the core 
capabilities the system must provide in order to successfully 
digitize diagrams and allow users to interact with them efficiently.

The input module is responsible for collecting diagram data from the user in various forms and preparing it for processing. To ensure flexibility, a multi-format image upload system is implemented that supports PNG, JPG, and WebP files with a drag-and-drop interface and file size validation. For algorithm-to-flowchart conversion, a multi-line text prompt area with character counting is provided. Additionally, the system enables direct JSON import with built-in validation to ensure compatibility with existing flow data structures.

The AI processing module manages communication with the Google Gemini API to perform robust shape recognition, text extraction, and connection detection. The system accurately identifies processes, decisions, start nodes, and input/output shapes from uploaded images, extracting their embedded text while preserving line breaks for multi-line labels. Furthermore, the module detects directional edges between nodes, assigning appropriate connection handles and estimating normalized coordinates to ensure consistent positioning. To ensure system reliability, a robust parsing mechanism with a regex-based fallback is implemented to handle AI response variability.

The editor module provides an interactive canvas with a rich set of manipulation features, including drag-and-drop functionality for adding shapes from a palette and sophisticated node resizing with text adaptation. For workflow safety, the editor implements snapshot-based undo and redo functionality with standard keyboard shortcuts. A centralized auto-layout engine, powered by the Dagre hierarchical algorithm, ensures clear alignment and optimal edge routing. Users can also interact with edges to add or edit labels, delete elements with automatic snapshot capture, and customize visual styles through color pickers and text formatting controls (bold, italic, and size adjustments).

The export module enables users to generate professional-quality outputs in multiple formats. PNG export is supported through the html-to-image library for easy sharing, while SVG export leverages React Flow’s built-in renderer for high-quality vector graphics. Furthermore, the system allows for JSON serialization of nodes and edges to ensure data portability and readability.

The storage module ensures long-term persistence by enabling users to save flowchart data with a UUID primary key and timestamp tracking in a PostgreSQL database. The module also supports retrieving flowcharts by UUID to reconstruct the editor state and provides update capabilities that automatically refresh the relevant timestamps upon modification.

\subsection{Non-Functional Requirements}
\label{sec:nonfunctional-req}

Non-functional requirements define the quality attributes of the 
system rather than its specific behaviors. They describe how well 
the system performs its functions, covering aspects such as speed, 
usability, reliability, and security, and are essential for ensuring 
a positive user experience and long-term system sustainability.

\subsubsection{Performance Requirements}

The system is designed to provide a responsive and efficient experience across all core operations. Image analysis is optimized to complete within 3 to 5 seconds for standard flowchart images, while text-to-flowchart generation is even faster, typically finishing within 2 to 3 seconds. Core editor interactions, including node manipulation, undo/redo, and auto-layout, are executed within 100 milliseconds to ensure a seamless user experience. To support high availability, the backend maintains API response times under 200 milliseconds for standard operations and achieves an initial page load time of under 2 seconds. The overall architecture is built to support at least 50 concurrent users without any noticeable performance degradation.

\subsubsection{Usability Requirements}

The platform prioritizes ease of use through an intuitive interface that requires minimal training, supported by clear, actionable error messages and helpful tooltips for all toolbar functions. A responsive design ensures comfortable usage across desktop, tablet, and mobile devices, while theme consistency between dark and light modes prevents visual fatigue. Additionally, the system supports standard keyboard shortcuts for critical operations like undo, redo, and deletion, further enhancing the productivity of experienced users.

\subsubsection{Reliability Requirements}

To ensure reliable operation, the system targets 99\% uptime during normal business hours and implements graceful error handling with robust fallback mechanisms for AI processing failures. Data integrity is maintained through full transaction support for all database operations, combined with rigorous input validation on both the client and server sides to prevent the processing of invalid data. Furthermore, the system supports comprehensive database backup and restore procedures to facilitate disaster recovery.

\subsubsection{Security Requirements}

Security is a primary consideration, with the Gemini API key protected through exclusive server-side storage in environment variables. All user inputs are sanitized to prevent injection attacks, and a strict CORS configuration restricts API access to verified origins. SQL injection protection is enforced through the use of parameterized queries via SQLAlchemy, while the maintenance of separate configuration files for development and production environments ensures a secure operational lifecycle.

%---------------------------------------------------------------
\section{System Architecture Overview}
\label{sec:arch-overview}
%---------------------------------------------------------------

% C4: Condensed Section 3.5 — detailed design moved to Chapter 4
% C5: Figure 3.1 reference removed to avoid duplication with Fig 4.1

The Automated Diagram Digitization platform is built on a three-tier 
architecture that separates the user interface, business logic, and 
data persistence into distinct, independently maintainable layers. 
This separation of concerns ensures that each layer can be developed, 
tested, and scaled independently without affecting the other 
components. A detailed description of the system architecture, 
including data flow diagrams, module design, API specifications, and 
database schema, is presented in Chapter~4.

%---------------------------------------------------------------
\section{System Modules}
\label{sec:system-modules}
%---------------------------------------------------------------

The Automated Diagram Digitization platform is organized into five 
functional modules, each responsible for a specific part of the 
overall workflow. This modular design simplifies development and 
maintenance while allowing individual components to be improved 
independently.

\subsection{Input Module}

The input module serves as the primary entry point of the system, 
accepting user-provided data in multiple forms and preparing it for 
further processing by the AI engine and the diagram editor. It 
handles three input methods: image upload with drag-and-drop 
interface, text prompt with multi-line textarea, and JSON import 
with file upload. It validates input formats, provides visual 
feedback, and prepares data for AI processing.

\subsection{AI Processing Module}

The AI processing module coordinates communication with the external 
Gemini AI service, managing both outgoing requests and the parsing 
of structured responses returned by the model. It interfaces with 
Google Gemini AI using custom system prompts for image analysis and 
text-to-flowchart conversion, parses JSON responses, handles errors 
with a regex-based fallback, and ensures proper ID generation for 
edges.

\subsection{Editor Module}

The editor module forms the interactive core of the platform, 
enabling users to view, manipulate, and refine the flowcharts 
generated by the AI processing module. It renders nodes and edges 
using React Flow with custom node components for each shape type, 
and implements drag-and-drop from sidebar, undo/redo with snapshot 
history, hierarchical auto-layout via Dagre, edge labeling via 
comprehensive styling options.

\subsection{Export Module}

The export module provides users with the ability to save and 
distribute their diagrams in formats suited to a range of 
professional and educational use cases. It provides PNG export using 
the \texttt{html-to-image} library, SVG export using React Flow's 
SVG renderer, and JSON export of the current flow state, with file 
downloads named according to the flowchart title.

\subsection{Storage Module}

The storage module ensures long-term persistence of diagram data, 
allowing users to save their work and retrieve it in future sessions 
without loss of layout or content information. It manages PostgreSQL 
database operations including save, retrieve, and update of 
flowcharts, implementing UUID-based identification, timestamp 
tracking, and JSONB storage for complete flow data.

%---------------------------------------------------------------
\section{User Roles and Permissions}
\label{sec:roles}
%---------------------------------------------------------------

In the current version of the Automated Diagram Digitization 
platform, the system operates without mandatory user authentication, 
making all features accessible to any user. A single general user 
role with full access to all features is implemented, as 
authentication is not required for basic usage. Future enhancements 
may include user authentication with role-based access control for 
collaborative features, as discussed in Chapter~8. The table below 
summarizes the permissions associated with the general user role. 
\textit{Refer Table~\ref{tab:roles}: User Roles and Access Levels.}

\begin{table}[ht]
\centering
\renewcommand{\arraystretch}{1.3}
\begin{tabular}{|l|p{10cm}|}
\hline
\textbf{Role} & \textbf{Permissions} \\
\hline
User (All) & Upload images, enter text prompts, import JSON \\
\hline
User (All) & Create, edit, and delete flowcharts \\
\hline
User (All) & Export diagrams (PNG, SVG, JSON) \\
\hline
User (All) & Toggle theme (dark/light mode) \\
\hline
User (All) & Save and retrieve flowcharts from the database \\
\hline
\end{tabular}
\caption{User Roles and Access Levels}
\label{tab:roles}
\end{table}

%---------------------------------------------------------------
\section{Conclusion}
\label{sec:ch3-conclusion}
%---------------------------------------------------------------

The system study has established clear and comprehensive requirements 
for the Automated Diagram Digitization project, addressing critical 
gaps in automated shape recognition, text extraction, connection 
detection, and interactive editing. The analysis of existing systems 
revealed significant inefficiencies in manual redrawing workflows, 
the absence of automation, and limited editing capabilities in 
current tools. These limitations highlight the need for a system 
that can intelligently interpret diagram structures while also 
providing flexible post-recognition editing capabilities.

The proposed system directly addresses these challenges through 
AI-powered recognition, multiple input methods, comprehensive editing 
features, and flexible export options. Feasibility analysis confirms 
that the solution is technically achievable using proven technologies 
such as FastAPI, React, PostgreSQL, and Google Gemini AI, which 
collectively provide a stable and scalable foundation.

From an operational perspective, the system offers a user-friendly 
interface that simplifies the overall workflow from diagram input to 
final export. Economically, the use of open-source technologies and 
cloud-based infrastructure reduces development and operational costs 
while providing long-term sustainability. Chapter~4 presents the 
detailed system design, translating the requirements established here 
into architectural, data flow, and module-level design 
specifications.